\section{Method and Approach}
\label{sec:data}

Our approach follows one practical rule: with limited resources, data quality and training stability are primary levers.
Therefore, we combine a compact model with strict preprocessing, synchronized augmentation, and modular feature fusion.

\subsection{Data preparation pipeline}

\paragraph{1) Unified dataset structure}
\textit{Guiding idea: keep training inputs consistent across sources.}
We prepare two different domains for training in one unified format with shared split/class structure.

\paragraph{2) Landmark metadata generation}
\textit{Guiding idea: use facial landmarks to guide the model toward important facial points and to better handle faces that are not perfectly aligned.}
For each sample, facial landmark coordinates are extracted, scaled to target resolution, and normalized.
The vectors are stored in \texttt{metadata.csv} so the model receives explicit spatial hints (e.g., around eyes, eyebrows, nose, and mouth) during geometry-appearance fusion.

\paragraph{3) Signal shaping}
\textit{Guiding idea: pass cleaner, more informative structure to the model earlier.}
Before model ingestion, images can be transformed with fixed edge operators (Sobel, Roberts, Laplacian, or raw baseline).
This creates a controlled preprocessing axis without changing classifier topology.

\subsection{Geometry-consistent augmentation}

\textit{Guiding idea: increase data diversity.}
When an image is transformed, landmark coordinates are transformed with the same geometric mapping.
Supported operations are random resized crop, rotation, translation, perspective warp, affine shear, and horizontal flip.
Perspective uses homography mapping in homogeneous coordinates; affine operations use explicit matrix transforms.
After transformation, coordinates are clipped and re-normalized.

\subsection{Imbalance-aware optimization}

\textit{Guiding idea: limit majority-class dominance in uneven datasets.}
Class imbalance is addressed with two modes: class-balanced undersampling and inverse-frequency loss weighting.
This balances gradient contributions across emotion classes without changing the core architecture.
